% PyMORGAN — Steady-state spectra chapter
\section{Steady-state spectra}\label{sec:steadyState}

The \code{pymorgan.steadyState} module provides lightweight wrappers for
absorption, emission and excitation spectra. Its design mirrors the
one-dimensional branch (chapter~\ref{sec:oneD}): data are held in a plain
object (\code{Spectrum}), file formats are resolved through a loader registry,
and presentation is driven by the shared \code{Settings}.

\subsection{Data model}

\subsubsection{\code{SpectrumKind}}

\code{SpectrumKind} is a string enum that labels the physical quantity
measured:

\begin{lstlisting}
class SpectrumKind(StrEnum):
    ABSORPTION = "absorption"
    EMISSION   = "emission"
    EXCITATION = "excitation"
\end{lstlisting}

\subsubsection{\code{Spectrum}}

A single steady-state measurement. The spectral axis \code{x} and signal
\code{y} are stored as \code{numpy} arrays; the remaining fields are metadata:

\begin{table}[h]
  \centering
  \begin{tabular}{lll}
    \toprule
    Attribute & Type & Meaning \\
    \midrule
    \code{x}       & \code{ndarray} & Spectral axis (values in \code{x\_units}) \\
    \code{y}       & \code{ndarray} & Signal (absorbance or intensity) \\
    \code{kind}    & \code{SpectrumKind} & Measurement type \\
    \code{x\_units}& \code{str} & Spectral-axis unit token (\code{"nm"}, \code{"cm-1"}, \code{"eV"}) \\
    \code{label}   & \code{str | None} & Legend label \\
    \code{source}  & \code{str | None} & File path (set automatically on load) \\
    \bottomrule
  \end{tabular}
  \caption{Attributes of a \code{Spectrum} object.}
  \label{tab:spectrum:fields}
\end{table}

\subsubsection{\code{SpectrumSeries}}

An ordered list of \code{Spectrum} objects for overlay plots. It supports
iteration, indexing and the standard sequence protocol. Construction from
multiple files returns a \code{SpectrumSeries} directly.

\subsection{Loading and the loader registry}

The module exposes a loader registry analogous to the 1-D and 2-D registries.
The default \code{"csv"} reader handles comma- and tab-delimited files with
automatic delimiter detection. Custom readers are registered with the
\code{@pm.register\_spectrum\_loader} decorator:\

\begin{lstlisting}
import pymorgan as pm

# List available loaders
print(pm.available_spectrum_loaders())

# Load a single spectrum
sp = pm.load_spectrum("sample.csv", kind="absorption")

# Specify units explicitly (overrides the reader's guess)
sp = pm.load_spectrum("sample.csv", kind="emission", x_units="nm", label="Sample A")

# Load directly via Dataset1D convenience method
sp2 = pm.Spectrum.from_file("sample.opus", kind="absorption", data_type="opus")
\end{lstlisting}

The currently registered loaders are \code{csv} (generic delimiter-detected
text), \code{opus} (Bruker OPUS binary, via \code{load.py}), \code{uv\_vis}
(Perkin-Elmer / Shimadzu UV-Vis text export) and \code{fluorimeter} (Horiba
FluorEssence \code{.fes} files).

A custom loader returns a three-element tuple:
\begin{lstlisting}
@pm.register_spectrum_loader("MySpectrometer")
def read_my_spectrometer(path):
    ...
    return x, y, x_units    # x_units may be None (falls back to "nm")
\end{lstlisting}

\subsection{Transformations}

Two per-\code{Spectrum} transformations return a new \code{Spectrum}
(the original is unchanged):

\begin{description}
  \item[\code{normalised()}] Scales the signal so that
    $\max(|y|) = 1$.
  \item[\code{to\_stimulated\_emission()}] Computes the stimulated-emission
    cross-section $F_\text{stim} \propto \bar{\nu}^4 \, F_\text{spont}$ (valid
    for emission spectra on a wavelength axis). Useful when overlaying on
    transient-absorption spectra.
\end{description}

\code{SpectrumSeries.normalised()} applies the same operation to every member
and returns a new series.

\subsection{Integration with the 1-D pipeline}

A \code{Spectrum} can be passed directly to \code{Dataset1D.plot\_spectra} as a
steady-state overlay:\

\begin{lstlisting}
ftir = pm.load_spectrum("background_FTIR.csv", kind="absorption")
data.plot_spectra([0.5, 1, 5], ss_overlay=ftir)
\end{lstlisting}

\code{as\_overlay\_dict()} converts a \code{Spectrum} to the
\code{\{"X": x, "Y": y\}} mapping that \code{plot\_spectra} also accepts,
for cases where you want to pass raw arrays.

\subsection{Plotting}

Both \code{Spectrum} and \code{SpectrumSeries} have a \code{plot()} method.
The axis-label convention and delimiter follow the active settings
(\S\ref{sec:settings}), and optional keyword arguments are forwarded to
\code{matplotlib.Axes.plot}.

\begin{lstlisting}
import pymorgan as pm

pm.load_settings("settings.toml")
pm.apply_style()

# --- Single spectrum ---
sp = pm.load_spectrum("sample.csv", kind="absorption", label="Sample")
sp.plot()

# --- Normalised emission ---
em = pm.load_spectrum("emission.csv", kind="emission", label="Emission")
em.normalised().plot()

# --- Series: multiple spectra coloured along a rainbow colourmap ---
series = pm.SpectrumSeries.from_files(
    ["a.csv", "b.csv", "c.csv"],
    kind="absorption",
    labels=["10 uM", "50 uM", "100 uM"],
)
series.plot(cmap="viridis", normalise=True)

pm.show_plots()
\end{lstlisting}

The per-call keyword arguments are identical to those of the 1-D plotters:
\code{ax=} to draw into an existing axis, \code{label\_style=} to override the
delimiter, and \code{normalise=} to scale the signal to unit peak.

